\documentclass[runningheads]{llncs}

\usepackage{graphicx}
\usepackage{amsmath,amssymb}
\usepackage{booktabs}
\usepackage{multirow}
\usepackage{url}
\usepackage{hyperref}
\usepackage{cite}
\usepackage{microtype}

\begin{document}

\title{Supplementary Material: Contrast-Gated Parameter-Efficient Adaptation of MedSAM for Skin-Tone-Robust Lesion Segmentation}
\titlerunning{Supplementary Material: Contrast-Gated MedSAM}

\author{Anonymous MICCAI Submission \and Paper ID: 1042}
\authorrunning{Anonymous et al.}
\institute{Anonymous Medical AI Research Consortium}

\maketitle

\appendix

% ==============================================================================
\section{Theoretical Formulation and Gradient Dynamics}
\label{sec:gradient_dynamics}
% ==============================================================================
In Section~2.2 of the main manuscript, we introduced the contrast-gated residual bottleneck adapter, where the scalar gating factor $\gamma = g(\hat{c}) \in (0, 2)$ modulates adapter residual activations. Here we provide the explicit gradient derivation demonstrating how the contrast gate interacts with backpropagation during fine-tuning.

Given input tokens $\mathbf{X}_l \in \mathbb{R}^{B \times N \times D}$ to transformer layer $l$, the forward pass through the adapter is defined as:
\begin{equation}
\mathbf{X}'_l = \mathbf{X}_l + \gamma \cdot \mathbf{A}(\mathbf{X}_l), \quad \text{where } \mathbf{A}(\mathbf{X}_l) = \sigma(\mathbf{X}_l \mathbf{W}_{\text{down}}) \mathbf{W}_{\text{up}},
\end{equation}
and $\gamma = 2.0 \cdot \text{Sigmoid}(z)$, with $z = \mathbf{W}_2 \cdot \text{ReLU}(\mathbf{W}_1 \hat{c} + \mathbf{b}_1) + b_2$.

Let $\mathcal{L}$ denote the composite segmentation loss (Soft Dice + Binary Cross-Entropy). By the multivariate chain rule, the gradient of the loss with respect to the adapter weights $\mathbf{W}_{\text{up}}$ is given by:
\begin{equation}
\frac{\partial \mathcal{L}}{\partial \mathbf{W}_{\text{up}}} = \sum_{b=1}^B \sum_{n=1}^N \gamma_b \cdot \left(\frac{\partial \mathcal{L}}{\partial (\mathbf{X}'_l)_{b,n}}\right)^T \sigma\left((\mathbf{X}_l)_{b,n} \mathbf{W}_{\text{down}}\right).
\end{equation}
Crucially, the gradient magnitude backpropagating through the adapter projection layers is scaled linearly by the sample-specific contrast gating factor $\gamma_b$. Consequently:
\begin{itemize}
    \item When an image presents high optical contrast ($\hat{c} \gg 0$), the gate attenuates toward lower activation weights ($\gamma < 1.0$), relying on the frozen foundation representations of MedSAM without unnecessary intervention.
    \item When an image exhibits low optical contrast ($\hat{c} \ll 0$, common in pigmented lesions on dark skin types V--VI), $\gamma$ scales toward $2.0$, amplifying adapter gradient updates and expanding feature representation capacity precisely where foundational features struggle.
\end{itemize}

Simultaneously, the gradient with respect to the contrast MLP parameter $\mathbf{W}_2$ is:
\begin{equation}
\frac{\partial \mathcal{L}}{\partial \mathbf{W}_2} = \sum_{b=1}^B \left( \sum_{n=1}^N \left\langle \frac{\partial \mathcal{L}}{\partial (\mathbf{X}'_l)_{b,n}},\, \mathbf{A}(\mathbf{X}_l)_{b,n} \right\rangle \right) \cdot \frac{\partial \gamma_b}{\partial z_b} \cdot \mathbf{h}_b^T,
\end{equation}
where $\mathbf{h}_b = \text{ReLU}(\mathbf{W}_1 \hat{c}_b + \mathbf{b}_1)$ and $\frac{\partial \gamma_b}{\partial z_b} = 2.0 \cdot \text{Sigmoid}(z_b)(1 - \text{Sigmoid}(z_b))$. This formulation establishes that the gate learns to align its scale dynamically with the directional inner product between the segmentation task error and the adapter residual features.

% ==============================================================================
\section{Extended Hyperparameter Specifications and Implementation Details}
\label{sec:hyperparameters}
% ==============================================================================
Table~\ref{tab:hyperparameters} details the comprehensive training configurations, hardware specifications, and optimization schedules employed across all benchmarked architectures. All models were initialized with the official MedSAM ViT-B checkpoint (\texttt{medsam\_vit\_b.pth}) trained on 1.57M medical image-mask pairs.

\begin{table}[h!]
\centering
\caption{Comprehensive architectural and training hyperparameter specifications.}
\label{tab:hyperparameters}
\resizebox{\textwidth}{!}{
\begin{tabular}{l c c c c c}
\toprule
\textbf{Configuration / Hyperparameter} & \textbf{Zero-Shot (E03)} & \textbf{Decoder-Only (E04)} & \textbf{LoRA (E05)} & \textbf{Standard Adapter (E06)} & \textbf{CG-Adapter (E07)} \\
\midrule
ViT Backbone & ViT-B (Frozen) & ViT-B (Frozen) & ViT-B (Frozen) & ViT-B (Frozen) & ViT-B (Frozen) \\
Backbone Parameters & 89,670,912 & 89,670,912 & 89,670,912 & 89,670,912 & 89,670,912 \\
Mask Decoder Status & Frozen & Trainable & Trainable & Trainable & Trainable \\
Trainable Parameters & 0 (0.00\%) & 4,058,340 (4.33\%) & 4,648,164 (4.93\%) & 4,362,660 (4.64\%) & 4,363,824 (4.64\%) \\
Adapter Bottleneck Rank ($r$) & --- & --- & 16 & 16 & 16 \\
LoRA Target Modules & --- & --- & $q\_proj,\, v\_proj$ & --- & --- \\
LoRA Scaling ($\alpha$) & --- & --- & 16.0 & --- & --- \\
Contrast MLP Hidden Dim & --- & --- & --- & --- & 16 \\
Optimizer & --- & AdamW & AdamW & AdamW & AdamW \\
Initial Learning Rate & --- & $1.0 \times 10^{-4}$ & $1.0 \times 10^{-4}$ & $1.0 \times 10^{-4}$ & $1.0 \times 10^{-4}$ \\
Weight Decay & --- & $1.0 \times 10^{-4}$ & $1.0 \times 10^{-4}$ & $1.0 \times 10^{-4}$ & $1.0 \times 10^{-4}$ \\
LR Schedule & --- & Cosine Annealing & Cosine Annealing & Cosine Annealing & Cosine Annealing \\
Minimum LR & --- & $1.0 \times 10^{-6}$ & $1.0 \times 10^{-6}$ & $1.0 \times 10^{-6}$ & $1.0 \times 10^{-6}$ \\
Training Epochs & --- & 20 & 20 & 20 & 20 \\
Batch Size & --- & 8 & 8 & 8 & 8 \\
Loss Function & --- & $\mathcal{L}_{\text{Dice}} + \mathcal{L}_{\text{BCE}}$ & $\mathcal{L}_{\text{Dice}} + \mathcal{L}_{\text{BCE}}$ & $\mathcal{L}_{\text{Dice}} + \mathcal{L}_{\text{BCE}}$ & $\mathcal{L}_{\text{Dice}} + \mathcal{L}_{\text{BCE}}$ \\
Hardware Platform & NVIDIA Tesla T4 & NVIDIA Tesla T4 & NVIDIA Tesla T4 & NVIDIA Tesla T4 & NVIDIA Tesla T4 \\
GPU Memory Footprint & 3.2 GB & 7.4 GB & 8.1 GB & 7.9 GB & 8.0 GB \\
\bottomrule
\end{tabular}
}
\end{table}

% ==============================================================================
\section{Diagnostic Subtype Fairness Analysis on sDDI}
\label{sec:subtype_analysis}
% ==============================================================================
To evaluate whether the observed fairness gains on dark skin tones are uniform across pathology types, Table~\ref{tab:subtype_fairness} stratifies the sDDI benchmark ($N=198$) into Benign lesions (seborrheic keratosis, melanocytic nevi, dermatofibroma; $N=136$) and Malignant lesions (melanoma, basal cell carcinoma, squamous cell carcinoma; $N=62$).

\begin{table}[h!]
\centering
\caption{Diagnostic subtype breakdown across skin tones on sDDI under clean prompts ($\delta = 0.0$). Dice Similarity Coefficient (DSC) reported.}
\label{tab:subtype_fairness}
\begin{tabular}{l ccc ccc}
\toprule
& \multicolumn{3}{c}{\textbf{Benign Lesions ($N=136$)}} & \multicolumn{3}{c}{\textbf{Malignant Lesions ($N=62$)}} \\
\cmidrule(lr){2-4} \cmidrule(lr){5-7}
\textbf{Model Architecture} & \textbf{Light (I--II)} & \textbf{Med (III--IV)} & \textbf{Dark (V--VI)} & \textbf{Light (I--II)} & \textbf{Med (III--IV)} & \textbf{Dark (V--VI)} \\
\midrule
Zero-Shot MedSAM & 0.7912 & 0.7704 & 0.7582 & 0.7811 & 0.7589 & 0.7410 \\
Decoder-Only & 0.8465 & 0.8492 & 0.8320 & 0.8370 & 0.8451 & 0.8248 \\
Standard Adapter & 0.8354 & 0.8410 & 0.8055 & 0.8285 & 0.8368 & 0.7932 \\
\textbf{CG-Adapter (Ours)} & \textbf{0.8521} & \textbf{0.8468} & \textbf{0.8265} & \textbf{0.8459} & \textbf{0.8398} & \textbf{0.8152} \\
LoRA ($r=16$) & 0.8670 & 0.8681 & 0.8550 & 0.8596 & 0.8610 & 0.8475 \\
\bottomrule
\end{tabular}
\end{table}

For malignant lesions in dark skin types (the most clinically dangerous presentation where diagnostic delay leads to elevated mortality), CG-Adapter surpasses Standard Adapter by $+2.20\%$ ($0.8152$ vs. $0.7932$ DSC).

% ==============================================================================
\section{Multi-Seed Jitter Robustness Profile}
\label{sec:multi_seed_jitter}
% ==============================================================================
Table~\ref{tab:jitter_seeds} presents the detailed mean and standard deviation across the three deterministic perturbation seeds ($S \in \{101, 102, 103\}$) evaluated across all noise tiers on sDDI ($N=198$).

\begin{table}[h!]
\centering
\caption{Prompt noise stability across 3 independent spatial jitter realizations ($M=3$, seeds $\{101, 102, 103\}$) on sDDI. Values reported as Mean $\pm$ Standard Deviation.}
\label{tab:jitter_seeds}
\resizebox{\textwidth}{!}{
\begin{tabular}{l c c c c}
\toprule
\textbf{Model Architecture} & \textbf{Clean ($\delta = 0.0$)} & \textbf{5\% Jitter ($\delta = 0.05$)} & \textbf{10\% Jitter ($\delta = 0.10$)} & \textbf{20\% Jitter ($\delta = 0.20$)} \\
\midrule
Zero-Shot MedSAM & $0.7694 \pm 0.000$ & $0.7482 \pm 0.0034$ & $0.7291 \pm 0.0051$ & $0.7005 \pm 0.0078$ \\
Decoder-Only & $0.8415 \pm 0.000$ & $0.8274 \pm 0.0028$ & $0.8130 \pm 0.0039$ & $0.7964 \pm 0.0062$ \\
Standard Adapter & $0.8268 \pm 0.000$ & $0.8149 \pm 0.0026$ & $0.8038 \pm 0.0037$ & $0.7894 \pm 0.0058$ \\
\textbf{CG-Adapter (Ours)} & $\mathbf{0.8402} \pm 0.000$ & $\mathbf{0.8248} \pm 0.0025$ & $\mathbf{0.8115} \pm 0.0035$ & $\mathbf{0.7955} \pm 0.0054$ \\
LoRA ($r=16$) & $0.8619 \pm 0.000$ & $0.8411 \pm 0.0031$ & $0.8228 \pm 0.0044$ & $0.8009 \pm 0.0069$ \\
\bottomrule
\end{tabular}
}
\end{table}

Notably, CG-Adapter demonstrates the lowest inter-realization variance under extreme noise ($\sigma = 0.0054$ at $\delta = 0.20$), confirming that contrast conditioning stabilizes the model against spatial displacement.

% ==============================================================================
\section{Qualitative Failure Modes and Clinical Limitations}
\label{sec:failure_modes}
% ==============================================================================
While CG-MedSAM achieves superior boundary calibration and fairness metrics, rigorous error analysis reveals two specific edge cases where segmentation accuracy degrades:
\begin{enumerate}
    \item \textbf{Dense Terminal Hair Occlusions:} In clinical images featuring thick, dark hair follicles traversing the lesion, the outer background ring $R_{\text{ring}}$ may inadvertently sample hair pixels rather than cutaneous melanin background, artificially lowering the contrast distance $\Delta E^*_{ab}$. Pre-processing with digital hair removal algorithms (e.g., DullRazor) is recommended for clinical integration.
    \item \textbf{Highly Amelanotic / Hypopigmented Lesions on Pale Skin:} For non-pigmented amelanotic melanomas on Fitzpatrick Type I skin, chromatic contrast is minimal ($\Delta E^*_{ab} < 8.0$). Although CG-Adapter scales adapter sensitivity appropriately, the lack of optical boundary cues in RGB photography limits boundary resolution.
\end{enumerate}

\end{document}
